% Generated draft: requires mathematical and source review.
\documentclass[11pt]{article}
\usepackage[margin=25mm]{geometry}
\usepackage{fontspec}
\setmainfont{Latin Modern Roman}
\usepackage{amsmath,amssymb,mathtools,unicode-math}
\setmathfont{Latin Modern Math}
\usepackage{enumitem,needspace,xurl,hyperref}
\setlength{\parindent}{0pt}
\setlength{\parskip}{5pt}
\setlength{\emergencystretch}{3em}
\newcommand{\N}{\mathbb{N}}
\newcommand{\Z}{\mathbb{Z}}
\newcommand{\Q}{\mathbb{Q}}
\newcommand{\R}{\mathbb{R}}
\newcommand{\C}{\mathbb{C}}
\newcommand{\F}{\mathbb{F}}
\newcommand{\A}{\mathbb{A}}
\newcommand{\PP}{\mathbb{P}}
\newcommand{\E}{\mathbb{E}}
\newcommand{\eps}{\varepsilon}
\newcommand{\dd}{\,\mathrm{d}}
\begin{document}
\section*{UnsolvedMath ID: 700430}
\textbf{Source problem code:} catalogue entry
\subsection{Definitions and mathematical statement}
Let \(G\) be a finite group. For a subgroup \(V\leq G\) and \(g\in G\), write \(V^g=g^{-1}Vg\), and define
\[
N_V(V^g)=\{v\in V:v^{-1}V^gv=V^g\}.
\]
A subgroup \(V\) has the trivial-intersection property (the \(TI\) property) if, for every \(g\in G\), either \(V^g=V\) or \(V\cap V^g=\{1\}\). The formulation below reads the source's term ``elementary \(TI\)-subgroup'' as an elementary abelian subgroup with this property. Thus, in the present problem, \(V\cong (\mathbb Z/2\mathbb Z)^n\) for some positive integer \(n\). The source does not explicitly expand the term ``elementary.''

Such a subgroup \(V\) is of \emph{non-root type} if
\[
\exists g_0\in G\qquad \{1\}\subsetneq N_V(V^{g_0})\subsetneq V.
\]
An involution is an element of order exactly \(2\). Write \([A,B]=1\) when every element of \(A\) commutes with every element of \(B\); in this case \(AB\) is a subgroup.

Describe all finite groups \(G\) for which there exists an elementary abelian \(TI\) \(2\)-subgroup \(V\leq G\) of non-root type satisfying
\[
\forall g\in G,\qquad
[V,V^g]=1\ \Longrightarrow\
\left(\forall t\in VV^g,\quad
(t\neq 1\ \land\ t^2=1)\ \Longrightarrow\
\exists x\in G\quad x^{-1}tx\in V\right).
\]
Here the final implication in the source is read as applying to every \(g\in G\), and conjugacy is read as conjugacy in \(G\). These scopes are implicit in the supplied wording and require confirmation against the cited source.
\subsection{Short English statement}
Describe the finite groups containing an elementary abelian \(TI\) \(2\)-subgroup \(V\) of non-root type such that, whenever \(V\) commutes with a conjugate \(V^g\), every involution in \(VV^g\) is conjugate in \(G\) to an element of \(V\). The expansion of ``elementary'' and the implicit quantifiers require source confirmation.
\subsection{Sources}
\begin{itemize}
\item Local UnsolvedMath registry: permanent ID \(700430\), source problem code \texttt{catalogue entry}, title ``Kourovka Notebook Problem 7.33.'' The supplied registry labels the problem open; this label is not an independently verified mathematical or literature assessment.
\item Supplied attribution: A. A. Makhn\"ev, proposed in 1980; \emph{The Kourovka Notebook (Unsolved Problems in Group Theory)}, 7th issue (1980), Problem 7.33, as cited in the supplied 21st-issue edition (2026), \texttt{arXiv:1401.0300v48}. The supplied collection description names E. I. Khukhro and V. D. Mazurov as editors.
\item The supplied registry record was created and updated on \texttt{2026-10-16T14:00:00Z}; its upstream revision is \texttt{d371ee917ac9647fbb06cd6fc529deec6f3dda82}. These are registry provenance details, not dates of independent source verification.
\end{itemize}
\end{document}
